\section{Evidence-Chain Forensics}
Let a trace be $T=(q,a,C,Z,M)$, where $q$ is a query, $a$ an answer,
$C=\{c_j\}$ retrieved contexts, $Z$ citation edges, and $M$ optional source and
execution metadata. Claim decomposition produces $A=\{a_i\}$. For each claim,
the verifier emits
\begin{equation}
 y_i=(v_i,e_i,z_i,s_i,r_i,f_i),
\end{equation}
where $v_i$ is a support verdict, $e_i$ a localized evidence span, $z_i$ a
citation state, $s_i$ a source-state assessment, $r_i$ a reason, and $f_i$ a
claim-level fix. An aggregate diagnosis $D(T)$ contains failure labels, one
primary observable root cause, an abstention decision, and a regression target.

Evidence-chain forensics imposes three invariants. \emph{Traceability}: a
supported verdict points to source text. \emph{Coverage honesty}: unavailable
outputs remain N/A rather than zero or inferred. \emph{Repair locality}: root
causes describe the earliest observable component supported by the trace, not
hidden implementation state. For example, missing evidence can support
``retrieval miss'' only when a supplied corpus or retrieval contract establishes
that the evidence should have been available; otherwise the safer label is
``should have abstained'' or ``corpus gap.''

Figure~\ref{fig:trace} illustrates a contradiction chain. The evidence says
refunds must be requested within 30 days while the answer claims 60 days. The
chain preserves the contradictory span and links it to answer overreach and a
specific answer constraint. This structure is intended to make diagnoses
inspectable by both developers and reviewers.

\begin{figure*}[t]
\centering
\begin{tikzpicture}[
  node distance=5mm and 5mm,
  box/.style={draw, rounded corners=2pt, align=center, minimum height=8mm, text width=22mm, font=\small},
  arrow/.style={-{Latex[length=2mm]}, thick}
]
\node[box] (trace) {RAG trace\\query, answer, contexts};
\node[box, right=of trace] (claims) {Claim\\decomposition};
\node[box, right=of claims] (evidence) {Evidence and\\citation checks};
\node[box, right=of evidence] (conditions) {Source and\\abstention checks};
\node[box, right=of conditions] (diagnosis) {Root cause and\\repair record};
\draw[arrow] (trace) -- (claims);
\draw[arrow] (claims) -- (evidence);
\draw[arrow] (evidence) -- (conditions);
\draw[arrow] (conditions) -- (diagnosis);
\end{tikzpicture}
\caption{ContextTrace production pipeline. Each stage emits a serializable
record used by downstream diagnosis and audit.}
\label{fig:pipeline}
\end{figure*}

\begin{figure}[t]
\centering
\begin{tikzpicture}[
  node distance=4mm,
  item/.style={draw, rounded corners=2pt, align=left, text width=68mm, font=\footnotesize, inner sep=4pt},
  arrow/.style={-{Latex[length=1.8mm]}, thick}
]
\node[item] (claim) {Claim: ``The policy allows a refund after 60 days.''};
\node[item, below=of claim] (span) {Evidence span: ``Refund requests must be made within 30 days.''};
\node[item, below=of span] (verdict) {Verdict: contradicted; citation target found};
\node[item, below=of verdict] (repair) {Diagnosis: answer overreach; repair: constrain the answer to the stated window};
\draw[arrow] (claim) -- (span);
\draw[arrow] (span) -- (verdict);
\draw[arrow] (verdict) -- (repair);
\end{tikzpicture}
\caption{Illustrative evidence-chain trace. The example describes the output
schema and is not a scored benchmark result.}
\label{fig:trace}
\end{figure}

